% Please do not change %
\documentclass[11pt]{article} % font size
\usepackage[margin=1in]{geometry} % margins
\usepackage{amsmath,amsthm,amssymb} % for the  common "math symbols"
\usepackage[shortlabels]{enumitem} % for enumeration
\usepackage{booktabs} % if you need tables
\usepackage{listings}
\usepackage{color}

\lstset{
  frame=none,
  xleftmargin=2pt,
  stepnumber=1,
  numbers=left,
  numbersep=5pt,
  numberstyle=\ttfamily\tiny\color[gray]{0.3},
  belowcaptionskip=\bigskipamount,
  captionpos=b,
  escapeinside={*'}{'*},
  language=haskell,
  tabsize=2,
  emphstyle={\bf},
  commentstyle=\it,
  stringstyle=\mdseries\rmfamily,
  showspaces=false,
  keywordstyle=\bfseries\rmfamily,
  columns=flexible,
  basicstyle=\small\sffamily,
  showstringspaces=false,
  morecomment=[l]\%,
}
% Feel free to include additional packages (if needed), but make sure it does not override the margin/font requirements.
%
\setlength{\parindent}{0pt} % no indent for new paragraphs.
\setlength{\parskip}{5pt} % distance between paragraphs, which will be used when you have a blank line in your LaTeX code.


%%%%%%%%%%%%%%%%%%%%%%%%%%%%%%%%%%%%%%%%%%%%%%%%%%%%%%%%%%%%%%%%%%%%%%%%%%%%%%
\title{Principles of Programming Languages: \\Problem Sheet 1 - First steps}
%
\author{}
\date{}
%
\begin{document}
%
\maketitle

%You may use the usual \LaTeX\ environments to structure your answers:
%\begin{align}
%        \label{basegnn}
%        \mathbf{h}_u^{(t+1)} = \sigma \bigg( \mathbf{W}_\text{self}^{(t)} \mathbf{h}_u^{(t)}  + \mathbf{W}_\text{neigh}^{(t)} \sum_{v \in N(u)} \mathbf{h}_v^{(t)} + \mathbf{b}^{(t)} \bigg). 
%\end{align}


%\begin{table}[h]
%\caption{A simple table.}
%\centering
%\begin{tabular}[t]{lrrrr}
%\toprule
%Graphs & Can & Be  & Fun & Too \\ 
%\midrule 
%$G_1$ & $1$  & $2$ & $3$ & $4$ \\ 
%$G_2$ & $-1$  & $-2$ & $-3$ & $-4$ \\ 
%$G_3$ & $0.1$  & $0.2$ & $0.3$ & $0.4$ \\ 
%\bottomrule
%\end{tabular}
%\label{tab}
%\end{table}

\section*{Question 1}

% To answer each part of the question, please simply use a paragraph command 
% \paragraph{(a)}
\lstinputlisting[language=Haskell, lastline=13]{../PS1.fun}

\clearpage
\section*{Question 2}


\paragraph{(a)}
\lstinputlisting[language=Haskell, lastline=13]{../PS1.hs}

\paragraph{(b)}
\lstinputlisting[language=Haskell, linerange={15-19}]{../PS1.fun}

\paragraph{(c)}

\lstinputlisting[language=Haskell, firstline=15]{../PS1.hs}

\section*{Question 3}
\begin{lstlisting}[language=Haskell]
Number 3
Variable "x"
Apply (Variable "+") (Variable "x", Number 3)
Apply (Variable "*") [Apply (Variable "+") (Variable "x", Number 3), Number 4]
(Let (Val "x" (Number 2)) (Apply (Variable "+") [Variable "x",Number 3]))
Let (Val "f" (Lambda ["x"] (Apply (Variable "+") [Variable "x",Number 3]))) (Apply (Variable "f") [Number 2])
Apply (Variable "f") [Apply (Variable "g") [Variable "x"]]
Apply [Apply (Variable "f") (Variable "g")] [Variable "x"]
\end{lstlisting}

\section*{Question 4}
Yes it is legal, value 5.
$x$ is the argument of the function $f$.
$g$ becomes a function that adds its argument to the argument of function $f$.


\section*{Question 5}
Yes it is legal, value 24.

\section*{Question 6}
$f$ is set as a function that adds 3 to its argument, 
and has nothing to do with $x$ except when referencing its value at assignment of $f$.

\section*{Question 7}
An Error: wrong number of args is raised.
With too little parameters, the value of unassigned parameters is unknown,
and the Error: "param" is not defined can get raised.
With too many parameters, although the extra parameters do not affect the result,
it can be extremely confusing to use.

It takes the value from the latest argument of the name
i.e. value of $x$ is from the third parameter, not the first.
It is sensible, but should be avoided.

\section*{Question 8}
Both expressions set the value of $x=e_1$ in $e_2$ - 
the first directly, the second, as an argument in a a lambda function.

\section*{Question 9}
\paragraph{(a)}
The expression is essentially syntactic sugar for
\begin{lstlisting}[language=Haskell]
(lambda (x1, x2, ..., xn) e)(e1, e2, ..., en)
\end{lstlisting}
Replace with the above construct in the parser.
% As an example, an abstract syntax could be
% \begin{lstlisting}
% LetAnd (
%   (Val "x1" (Variable "e1")), 
%   (Val "x2" (Variable "e2")), ..., 
%   (Val "xn" (Variable "en"))
% ) (Variable "e")
% \end{lstlisting}
% Then define \texttt{eval} for \texttt{LetAnd}
% \begin{lstlisting}[language=Haskell]
%   eval (LetAnd ds e1) env =
%     eval (Apply (Lambda (fst xes) e1) (snd xes)) env
%       where xes = unzip ds
% \end{lstlisting}

\paragraph{(b)}
Replace from the parser the expression with the construct it is syntactic sugar for.
\paragraph{(c)}
No idea.

\section*{Question 10}
Similar to how round brackets () are recognised, with an expression between them,
square brackets [] can too.
When square brackets are used, the expressions in between, separated by commas,
can be used as arguments to the \texttt{list} primitive under \texttt{eval}.

\end{document}